% Generated by tools/textbook/export_listings.js from tools/decks/mastery.js and tools/concepts/registry.json: do not edit.
\begin{chapstatus}{Stato di apprendimento di questo capitolo}
\masterychip{now}{DA COMPRENDERE OGGI}\enspace la virgola fissa; complemento a uno ed eccesso k; la virgola mobile IEEE 754; il nome RV32IMAC.\par\smallskip
\masterychip{use}{DA UTILIZZARE OGGI}\enspace il programma completo, gli strumenti di PRG32 (si usa, non si studia, in questo corso).\par\smallskip
\masterychip{later}{DA COMPRENDERE PIÙ AVANTI}\enspace niente.
\end{chapstatus}
\begin{chapstatus}{Percorso del capitolo}
Alla fine del capitolo saprete:
\begin{enumerate}[nosep,leftmargin=2em]
\item rappresentare un numero reale in virgola fissa e calcolarci con le istruzioni intere
\item decodificare un numero IEEE 754 in precisione singola
\item leggere il nome RV32IMAC e dire che cosa il kit sa e non sa calcolare in hardware
\end{enumerate}
Si appoggia su: solo l'installazione del Capitolo~\ref{ch:setup}. Prossimo passo: Capitolo~\ref{ch:verm1p1}.
\end{chapstatus}
